\documentclass[conference]{IEEEtran}
\IEEEoverridecommandlockouts

\usepackage{cite}
\usepackage{amsmath,amssymb,amsfonts}
\usepackage{algorithmic}
\usepackage{graphicx}
\usepackage{textcomp}
\usepackage{xcolor}
\usepackage{booktabs}
\usepackage{url}
\usepackage{listings}

\lstset{
  basicstyle=\ttfamily\scriptsize,
  breaklines=true,
  frame=single,
  columns=fullflexible
}

\def\BibTeX{{\rm B\kern-.05em{\sc i\kern-.025em b}\kern-.08em
    T\kern-.1667em\lower.7ex\hbox{E}\kern-.125emX}}

\begin{document}

\title{Cognitive-Access: Automated Cognitive Accessibility and Human Value Inspection for AI Coding Assistants\\
\thanks{This research was supported by Deakin University, School of Information Technology.}
}

\author{
\IEEEauthorblockN{Harjot Chauhan}
\IEEEauthorblockA{\textit{School of Information Technology} \\
\textit{Deakin University}\\
Geelong, Australia \\
harjot@deakin.edu.au}
\and
\IEEEauthorblockN{Davoud Mougouei}
\IEEEauthorblockA{\textit{School of Information Technology} \\
\textit{Deakin University}\\
Geelong, Australia \\
davoud.mougouei@deakin.edu.au}
}

\maketitle

\begin{abstract}
Conversational and agentic AI coding assistants (e.g., GitHub Copilot, Cursor) are transforming modern software development. However, existing automated quality evaluations focus almost exclusively on code correctness, syntax validation, and security vulnerabilities, overlooking the cognitive accessibility of human-AI interactions. For neurodivergent software engineers—particularly those with ADHD, anxiety, or executive dysfunction—AI assistants frequently induce severe cognitive friction by outputting unstructured code dumps, unannounced session expiries, and irreversible bulk deletions without safety confirmations. In this demonstration, we present \textsc{Cognitive-Access}, an automated testing and inspection tool for AI coding assistant interactions. \textsc{Cognitive-Access} introduces an empirically grounded 3-tier traceability framework linking international standards (W3C COGA, W3C WCAG 2.2, EU AI Act) directly to cognitive-science mechanisms (working memory capacity, ambiguity intolerance, behavioral inhibition) and Schwartz Human Values. The tool operates via both a lightweight CI/CD command-line interface and an interactive Web Inspector, utilizing deterministic static analyzers with an optional zero-cost Google Gemini API judge. We demonstrate its utility on authentic developer interaction transcripts from Stack Overflow and open-source repositories.
\end{abstract}

\begin{IEEEkeywords}
Cognitive Accessibility, AI Coding Assistants, Human Values, Neurodiversity, Software Testing, ICSE Demonstration.
\end{IEEEkeywords}

\section{Introduction and Motivation}
AI coding assistants powered by Large Language Models (LLMs) are now ubiquitous across software engineering. While these tools accelerate boilerplate generation and refactoring, their human-computer interaction (HCI) affordances often conflict with the cognitive needs of neurodivergent software engineers \cite{barkley1997, carleton2006, cowan2001}.

Traditional software accessibility focuses primarily on perceptual and motor accommodations (e.g., screen reader compatibilities, color contrast under WCAG 2.1). In contrast, \textit{cognitive accessibility} addresses how software systems interact with human executive functions, including working memory limits, uncertainty tolerance, attention switching, and impulse inhibition \cite{perera2019, steel2007}. When an assistant outputs 40 lines of unstructured code without micro-steps, it exceeds the classical human working memory limit of $4 \pm 1$ conceptual chunks \cite{cowan2001}. When an assistant initiates bulk code replacements without an undo checkpoint, it strains developer inhibition mechanisms \cite{barkley1997}. Furthermore, when an assistant presents opaque modifications without explaining its underlying rationale, it triggers anxiety and ambiguity intolerance \cite{carleton2006}, compromising the developer's human values of \textit{Self-Direction} and \textit{Personal Security} \cite{schwartz2012, perera2019}.

To bridge this gap, we present \textsc{Cognitive-Access}, an open-source testing tool designed to evaluate AI coding assistant interactions against cognitive accessibility standards and human values. The tool provides:
\begin{enumerate}
    \item A rigorous 3-tier mapping framework linking regulatory standards directly to empirical cognitive science literature and Schwartz Human Values.
    \item A hybrid execution engine combining offline deterministic static analysis with zero-cost LLM-judged heuristics via the Google Gemini API.
    \item An interactive Web Inspector Dashboard and a CI/CD-ready CLI that generate audit reports with dynamic sentence-level quote provenance.
\end{enumerate}

\section{The 3-Tier Traceability Framework}
A core limitation of conventional accessibility checkers is the semantic gap between high-level guidelines and underlying developer needs. \textsc{Cognitive-Access} addresses this through a formal 3-tier traceability architecture (Fig.~\ref{fig:traceability}).

\subsection{Tier 1: Standards and Regulations}
Tier 1 synthesizes actionable requirements from international regulatory and accessibility standards:
\begin{itemize}
    \item \textbf{W3C COGA:} Making Content Usable for People with Cognitive and Learning Disabilities, covering Pattern 4.2.1 (Make Purpose Clear), Pattern 4.2.4 (Make Each Step Clear), and Pattern 4.5.2 (Support Undo/Go Back).
    \item \textbf{W3C WCAG 2.2:} Success Criterion 3.3.4 (Error Prevention for Data Deletion) and SC 2.2.6 (Timeouts Warning).
    \item \textbf{EU AI Act (Reg. 2024/1689):} Article 13 (Transparency and Rationale Provision) and Article 14 (Human Oversight Checkpoints).
\end{itemize}

\subsection{Tier 2: Cognitive-Science Mechanisms}
Rather than treating standards as arbitrary rules, Tier 2 grounds each requirement in peer-reviewed empirical psychology. Every rule in \textsc{Cognitive-Access} resolves to an exact sentence in an authoritative research corpus verified from original publisher PDFs:
\begin{itemize}
    \item \textit{Working Memory Overload:} Grounded in Cowan (2001) \cite{cowan2001}, which demonstrates that human central capacity is limited to $4 \pm 1$ chunks in the absence of chunking or rehearsal.
    \item \textit{Ambiguity Intolerance:} Grounded in Carleton et al. (2006) \cite{carleton2006}, defining Intolerance of Uncertainty as the tendency to perceive ambiguous situations as threatening.
    \item \textit{Behavioral Inhibition Deficit:} Grounded in Barkley (1997) \cite{barkley1997}, establishing that executive function deficits in ADHD impair the ability to delay prepotent responses.
    \item \textit{Time Blindness:} Grounded in Noreika et al. (2012) \cite{noreika2012} and Smith et al. (2002) \cite{smith2002}, demonstrating objective temporal estimation deficits in neurodivergent individuals.
\end{itemize}

\subsection{Tier 3: Schwartz Human Values}
Tier 3 maps cognitive dysfunctions to Schwartz's Universal Human Values Theory as adapted to Software Engineering by Perera et al. \cite{perera2019, schwartz2012}:
\begin{itemize}
    \item \textit{Security---Personal:} Safety and cognitive stability in one's immediate working environment.
    \item \textit{Self-Direction---Thought:} Freedom to cultivate one's own ideas and understand code modifications independently.
    \item \textit{Self-Direction---Action:} Freedom to determine one's actions and retain autonomous oversight over codebase modifications.
\end{itemize}

\begin{table}[t]
\caption{Representative 3-Tier Traceability Mapping in Cognitive-Access}
\label{tab:mapping}
\centering
\scriptsize
\begin{tabular}{p{1.8cm}p{2.3cm}p{1.6cm}p{1.4cm}}
\toprule
\textbf{Tier 1 Standard} & \textbf{Tier 2 Cognitive Concept} & \textbf{Verified Source} & \textbf{Tier 3 Value} \\
\midrule
W3C COGA 4.2.4 & Working Memory Overload & Cowan (2001) & Security---Personal \\
W3C COGA 4.2.1 & Ambiguity Intolerance & Carleton (2006) & Security---Personal \\
WCAG 2.2 SC 3.3.4 & Inhibition Deficit & Barkley (1997) & Security---Personal \\
WCAG 2.2 SC 2.2.6 & Time Blindness & Noreika (2012) & Security---Personal \\
EU AI Act Art. 13 & Ambiguity Intolerance & Carleton (2006) & Self-Direction---Thought \\
EU AI Act Art. 14 & Planning Deficit & Steel (2007) & Self-Direction---Action \\
\bottomrule
\end{tabular}
\end{table}

\section{System Architecture and Implementation}
\textsc{Cognitive-Access} is implemented in Python 3.12 with modular extensibility for both command-line workflows and interactive browser inspection.

\subsection{Execution Modes and API Configuration}
The tool employs a hybrid execution architecture:
\begin{enumerate}
    \item \textbf{Deterministic Mode (Default Offline):} Evaluates interaction text and code diffs using Abstract Syntax Tree (AST) pattern matching, regular expressions, and structural heuristics. It detects bulk deletions ($>15$ lines removed without confirmation), unannounced timeouts, and missing micro-step numbering. This mode requires zero API keys and functions completely offline.
    \item \textbf{LLM-Judged Mode (Free-Tier Gemini API):} When inferential reasoning is required (e.g., evaluating whether an explanation adequately resolves technical ambiguity or plain-language clarity), the tool queries Google's free-tier Gemini API (\texttt{gemini-1.5-flash}). The API key is loaded automatically from an environment variable (\texttt{GEMINI\_API\_KEY}) or a local \texttt{.env} file. If no key is detected, the runner gracefully falls back to deterministic-only mode without crashing.
\end{enumerate}

\subsection{Interactive Web Inspector}
In addition to the CLI, \textsc{Cognitive-Access} provides an interactive, client-side Web Inspector Dashboard (HTML5, Tailwind CSS, JavaScript). Developers can paste live Copilot/Cursor transcripts, adjust inspection parameters, inspect visual 3-tier traceability cards, and view human-value concern gauges in real time.

\begin{figure}[t]
\centering
\fbox{\begin{minipage}{0.95\linewidth}
\ttfamily\scriptsize
\$ python -m tool.cli test transcripts/copilot\_session.txt\\
================================================================================\\
COGNITIVE-ACCESS TESTING TOOL REPORT (ICSE 2027 DEMO TRACK)\\
================================================================================\\
Target File: transcripts/copilot\_session.txt\\
Execution Mode: ALL | Gemini LLM Judge: ACTIVE (Free Tier)\\
Total Violations: 2 (High: 1, Medium: 1)\\[2pt]
[1] COGA-4.2.4-WM [DETERMINISTIC] (HIGH) --- Unstructured Output\\
\hspace*{1em}Tier 1: W3C COGA Section 4.2 Pattern 4.2.4 (Make Each Step Clear)\\
\hspace*{1em}Tier 2: Working Memory Overload (Cowan 2001)\\
\hspace*{1em}Quote: "The present target article brings together a wide variety..."\\
\hspace*{1em}Tier 3: Security---Personal ("Safety in one's immediate environment.")\\
--------------------------------------------------------------------------------\\
[2] EU-AI-ACT-ART13 [LLM\_JUDGED] (MEDIUM) --- Opaque AI Output\\
\hspace*{1em}Tier 1: EU AI Act Article 13 (Transparency and Information)\\
\hspace*{1em}Tier 2: Ambiguity Intolerance (Carleton 2006)\\
\hspace*{1em}Quote: "Popular alternative measures are either longer than the IUS..."\\
\hspace*{1em}Tier 3: Self-Direction---Thought ("Freedom to cultivate ideas...")
\end{minipage}}
\caption{Terminal output from \textsc{Cognitive-Access} showing dynamic 3-tier diagnostic report.}
\label{fig:cli_output}
\end{figure}

\section{Empirical Evaluation and Benchmark}
To assess \textsc{Cognitive-Access} without circularity or synthetic bias, we constructed an authentic evaluation benchmark comprising 20 real-world developer-assistant interaction transcripts sourced from Stack Overflow and open-source discussion threads. Each transcript represents an authentic development task (e.g., DirectX shader debugging, Electron sandboxing, Avalonia UI refactoring).

Ground-truth labels were established through independent human adjudication across the 6 core criteria. We evaluated the tool in hybrid mode (Deterministic + Gemini Free Tier) against this ground truth.

\begin{table}[h]
\caption{Evaluation Results on 20 Authentic Developer Transcripts}
\label{tab:eval}
\centering
\scriptsize
\begin{tabular}{lcccccc}
\toprule
\textbf{Criterion Rule ID} & \textbf{TP} & \textbf{FP} & \textbf{FN} & \textbf{Precision} & \textbf{Recall} & \textbf{F1} \\
\midrule
COGA-4.2.4-WM & 8 & 0 & 0 & 1.00 & 1.00 & 1.00 \\
COGA-4.2.1-AMBIG & 11 & 1 & 0 & 0.92 & 1.00 & 0.96 \\
WCAG-3.3.4-INHIB & 3 & 0 & 0 & 1.00 & 1.00 & 1.00 \\
WCAG-2.2.6-TIME & 2 & 0 & 0 & 1.00 & 1.00 & 1.00 \\
EU-AI-ACT-ART13 & 11 & 0 & 0 & 1.00 & 1.00 & 1.00 \\
EU-AI-ACT-ART14 & 5 & 0 & 0 & 1.00 & 1.00 & 1.00 \\
\midrule
\textbf{Overall Benchmark} & \textbf{40} & \textbf{1} & \textbf{0} & \textbf{0.98} & \textbf{1.00} & \textbf{0.99} \\
\bottomrule
\end{tabular}
\end{table}

\subsection{Failure Analysis and Boundary Cases}
The evaluation revealed one instructive false positive on \texttt{COGA-4.2.1-AMBIG} (Precision = 0.92). In this instance, a developer pasted a condensed stack trace snippet where the assistant immediately responded with a direct patch. While the response lacked an explicit introductory purpose statement (triggering the heuristic), human adjudication noted that the immediate context rendered the goal obvious to the user. This demonstrates the necessity of hybrid architectures where developers can inspect and calibrate rule sensitivities.

\section{Demonstration Walkthrough and Video}
During the demonstration session, attendees will experience:
\begin{enumerate}
    \item \textbf{Live Interaction Audit:} Running \textsc{Cognitive-Access} against real Copilot and Cursor chat sessions in the terminal and viewing real-time 3-tier diagnostics.
    \item \textbf{Web Inspector Exploration:} Interactively testing code diffs in the browser and observing how unbuffered deletions trigger the Barkley (1997) inhibition alert.
    \item \textbf{Remediation Guidance:} Observing how restructuring output into numbered micro-steps with explicit rationales eliminates violations and aligns with Schwartz Human Values.
\end{enumerate}
A recorded 4-minute demonstration video and reproducible codebase are available on GitHub.\footnote{\url{https://github.com/1HarryChauhan/sel703-placement}}

\section{Conclusion}
\textsc{Cognitive-Access} pioneers automated cognitive accessibility testing for AI coding assistants. By bridging regulatory standards with empirical cognitive science and human values, it provides software engineering teams with an automated, transparent, and reproducible mechanism to ensure AI coding tools support all developers equitably.

\begin{thebibliography}{00}
\bibitem{barkley1997} R. A. Barkley, ``Behavioral inhibition, sustained attention, and executive functions: Constructing a unifying theory of ADHD,'' \textit{Psychological Bulletin}, vol. 121, no. 1, pp. 65--94, 1997.
\bibitem{carleton2006} R. N. Carleton, M. A. P. Norton, and G. J. G. Asmundson, ``Fearing the unknown: A short version of the Intolerance of Uncertainty Scale,'' \textit{Journal of Anxiety Disorders}, vol. 21, no. 1, pp. 105--117, 2007.
\bibitem{cowan2001} N. Cowan, ``The magical number 4 in short-term memory: A reconsideration of mental storage capacity,'' \textit{Behavioral and Brain Sciences}, vol. 24, no. 1, pp. 87--114, 2001.
\bibitem{noreika2012} V. Noreika, C. M. Falter, and K. Rubia, ``Timing deficits in attention-deficit/hyperactivity disorder (ADHD): Evidence from neurocognitive and neuroimaging studies,'' \textit{Neuropsychologia}, vol. 51, no. 2, pp. 235--266, 2013.
\bibitem{perera2019} H. Perera et al., ``A Study on the Matching between Schwartz Human Values and Software Engineering Values,'' in \textit{Proc. IEEE Int. Requirements Eng. Conf. (RE)}, 2019, pp. 104--114.
\bibitem{schwartz2012} S. H. Schwartz et al., ``Refining the theory of basic individual values,'' \textit{Journal of Personality and Social Psychology}, vol. 103, no. 4, pp. 663--688, 2012.
\bibitem{smith2002} A. B. Smith, K. Rubia, and E. Taylor, ``Evidence of pure time perception deficits in children with ADHD,'' \textit{Journal of Child Psychology and Psychiatry}, vol. 43, no. 4, pp. 529--542, 2002.
\bibitem{steel2007} P. Steel, ``The nature of procrastination: A meta-analytic and theoretical review of quintessential self-regulatory failure,'' \textit{Psychological Bulletin}, vol. 133, no. 1, pp. 65--94, 2007.
\end{thebibliography}

\end{document}
